\chapter*{\centerline{Abstract}}
\markboth{\MakeUppercase{Abstract}}{}
\iftoggle{fulltoc}{
  \addcontentsline{toc}{chapter}{Abstract}
}{}

Musculoskeletal humanoid robots enable controlled experiments in biological movement and motivate compliant locomotion in unstructured environments. Their development requires reliable actuator and joint models, physiological evidence, and tools for testing neural controllers with different mechanics. I address these needs through four connected contributions. First, characterization of Festo braided pneumatic actuator (BPA) across resting lengths, contractions, and pressures produced a resting-length-dependent maximum-force model and a nondimensional three-coefficient  force surface. Joint-torque design factors then account for length offsets, effective series compliance, and wrapping-induced loss of usable actuator length. Incorporating these factors into BPA routing redesign enabled paired knee actuators to exceed selected human-muscle reference torques at measured configurations. I created the Sensory Afferent Database to connect experimental findings and model assumptions through searchable notes and citation navigation. Graphical neural blocks and simulation interfaces support controller transfer and testing on highly actuated musculoskeletal models, including Gait2392, in MuJoCo and Simulink/Simscape. Representative evaluations establish circuit responses, implementation agreement, and virtual-harness-assisted balance and walking behaviors. These contributions connect experimentally informed robot design with neuromechanical experiments and potential prosthetic and orthotic applications. 